\documentclass[11pt]{article}

\usepackage[review]{acl}   % change to \usepackage{acl} for camera-ready

\usepackage{times}
\usepackage{latexsym}
\usepackage[T1]{fontenc}
\usepackage[utf8]{inputenc}
\usepackage{microtype}
\usepackage{inconsolata}
\usepackage{graphicx}
\usepackage{booktabs}
\usepackage{amsmath,amssymb}

% ── Title ──────────────────────────────────────────────────────────────────────
\title{TITLE_GOES_HERE}

% Single author:
%   \author{Name \\ Affiliation \\ \texttt{email@domain}}
%
% Multiple authors (same institution):
%   \author{Author 1 \and Author 2 \\ Affiliation \\ \texttt{email}}
%
% Multiple authors (different institutions) — use \And / \AND:
\author{First Author \\
  Affiliation / Address \\
  \texttt{email@domain}
\And
  Second Author \\
  Affiliation / Address \\
  \texttt{email@domain}
}

% ── Document ───────────────────────────────────────────────────────────────────
\begin{document}
\maketitle

\begin{abstract}
% 1–2 paragraphs: problem statement, approach, key results, and significance.
% ACL abstracts are typically 150–250 words.
\end{abstract}

\section{Introduction}
% Motivate the problem and state the research question.
% Summarise the contributions (bold bullets are common in ACL papers).
% End with a one-sentence roadmap of the remaining sections.

\section{Related Work}
% Survey prior work grouped thematically or by approach.
% Explain how this paper differs from or extends each group.

\section{Methodology}
% Describe the approach in enough detail for a reader to replicate it.
% Include task formulation, model architecture, algorithms, and design choices.

\section{Experiments}
% Describe datasets, baselines, evaluation metrics, and hyperparameters.
% State any data splits, preprocessing steps, or annotation procedures.

\section{Results}
% Report the main quantitative findings (tables and figures are encouraged).
% Highlight statistical significance and ablation outcomes where applicable.

\section{Discussion}
% Interpret the results: what do they mean beyond the numbers?
% Address unexpected findings, failure modes, and scope of the conclusions.

\section{Conclusion}
% Summarise the key contributions and findings in 1–2 paragraphs.
% State implications for the field and suggest concrete directions for future work.

\section{Limitations}
% Required by ACL venues.
% Be specific: data coverage, model scale, language/domain scope, compute, etc.

% ── Bibliography (must appear before \appendix in ACL format) ──────────────────
\bibliography{references}

% ── Appendix ───────────────────────────────────────────────────────────────────
\appendix

\section{Appendix}
\label{sec:appendix}
% Supplementary material: proofs, extended tables, additional experiments, etc.
% Appendix sections are numbered A, B, … automatically after \appendix.

\end{document}
